\section{Context, scope, and conventions}
\label{dhr:sec:context}

\subsection{The continuation problem}
The canonical \repo\ manuscript already has a substantial Gamma--harmonic
and Hurwitz--Stieltjes calculus. In particular, the inspected mixed-Dougall
chapter proves a normally convergent two-parameter Gamma germ, all its
fixed mixed derivatives, a weighted trigamma square, and a quadratic
zeta-tail family \cite{repo:readme,repo:dougall}. The manuscript also records
an all-order Gauss proof of the corrected pure and mixed Choi harmonic
families. Re-proving the one-variable coefficient engine would not by
itself be a meaningful continuation.

We instead ask what happens when a squared Vandermonde couples an arbitrary
number of summation indices. There are two distinct questions. First,
can the Gamma product still supply a correctly normalized Stieltjes
subtraction when the sum ceases to converge? Second, does the increase
in the number of indices force a comparable increase in polylogarithmic
depth? We answer both for the first resonant determinantal family.

The first answer requires attention to a distinction easily hidden by
formal Laurent expansion: a spectral finite part is generally \emph{not}
the constant term of a hard cutoff. The residue amplitude varies with
the spectral parameter. Its derivative is precisely the discrepancy.
The second answer is more economical than the number of variables
suggests: for simple-pole rational weights at critical decay, only the
bottom-right moment contains a harmonic number. This gives depth at most
two at every particle number.

\subsection{Main results and what they do not claim}
Theorem~\ref{dhr:thm:gauss} evaluates the convergent discrete partition
function by elementary determinant algebra and Gauss summation. Its
classical ingredients are not presented as new. Theorems
\ref{dhr:thm:resonance} and \ref{dhr:thm:jets} give the normally convergent
resonance calculus and its all-order Stieltjes consequences. Theorems
\ref{dhr:thm:rank-one} and \ref{dhr:thm:closure} give the finite-cutoff
rank-one harmonic theorem and its all-particle depth-two consequence.
The concrete formulas in Section~\ref{dhr:sec:explicit} are exact
applications, not identities inferred from a numerical relation search.

These are results proved in the present report. An exhaustive literature
comparison has not established their global novelty. Nor do they prove
a lower bound on the depth of any actual period. The manuscript explicitly
distinguishes formal cyclotomic quotients from numerical period
independence \cite{repo:quotients}; we retain that distinction. The short
$S_6$ and $S_8$ identities remain outside the results proved here.

\subsection{Notation and analytic conventions}
For $k\geq0$, write $(x)_k=\Gamma(x+k)/\Gamma(x)$ for a rising
factorial, and $\fall{x}{k}=x(x-1)\cdots(x-k+1)$ for a falling factorial.
Empty products and empty determinants are one. We use
\[
 H_n^{(s)}=\sum_{j=1}^n j^{-s},\qquad H_n=H_n^{(1)},\qquad H_0^{(s)}=0.
\]
The gamma logarithmic derivative is $\psi=\Gamma'/\Gamma$. Our
Stieltjes convention is
\begin{equation}\label{dhr:eq:stieltjes-convention}
 \zeta(1+u,\alpha)=\frac1u+
 \sum_{k\geq0}\frac{(-1)^k}{k!}\gamma_k(\alpha)u^k,
 \qquad \Rea\alpha>0.
\end{equation}
Thus $\gamma_0(\alpha)=-\psi(\alpha)$ and $\gamma_0(1)=\gamma$.
Derivatives $\zeta^{(k)}(s,\alpha)$ are in the first, or order,
variable unless specified otherwise. We use the standard branches of
Gamma and Hurwitz zeta on the stated right half-planes
\cite{dlmf:gamma,dlmf:hurwitz,dlmf:polygamma}.

Within $|z|<1$, define
\begin{equation}\label{dhr:eq:polylogs}
 \Li_s(z)=\sum_{n\geq1}\frac{z^n}{n^s},\qquad
 \Li_{p,1}(z,1)=\sum_{n>m\geq1}\frac{z^n}{n^pm}.
\end{equation}
The latter convention agrees with the two-argument convention
$\Li_{p,1}(x,y)=\sum_{n>m\geq1}x^ny^m/(n^pm)$.
In particular $\Li_{1,1}(z,1)=\tfrac12\log^2(1-z)$. All formulas
initially in the disk use the logarithm analytic there and zero at
$z=0$. Rational Laurent factors in $z$ in the displayed reductions
have removable singularities after the full expression is assembled.
Analytic continuation means continuation of this common germ, not an
independent choice of branch for each summand.

The complete exponential Bell polynomials are normalized by
\begin{equation}\label{dhr:eq:bell}
 \exp\left(\sum_{j\geq1}x_j\frac{t^j}{j!}\right)
 =\sum_{k\geq0}\Bell_k(x_1,\ldots,x_k)\frac{t^k}{k!}.
\end{equation}
For example, $\Bell_1=x_1$, $\Bell_2=x_1^2+x_2$, and
$\Bell_3=x_1^3+3x_1x_2+x_3$.

Whenever a series of subtracted brackets is asserted to converge,
the brackets are the summands. This does not license splitting the
series into separately divergent components. All index and parameter
derivatives in the results below are fixed finite orders.
